\section*{Hoofdstuk \ref{ch:g11:electric-gravitational-fields} --- Elektrische en gravitatievelden}

\begin{solution}{exo:g11:electric-gravitational-fields:1}
$E = \num{5.0e-3} / \num{2.0e-6} = \qty{2.5e3}{N/C}$, naar het noorden.
Op $q'$: $F = \num{4.0e-6} \times \num{2.5e3} = \qty{1.0e-2}{N}$, naar
het \emph{zuiden} ($q' < 0$).
\end{solution}

\begin{solution}{exo:g11:electric-gravitational-fields:2}
\emph{1.} $E = \num{8.99e9} \times \num{5.0e-9} / 0.30^2
\approx \qty{5.0e2}{N/C}$, van $O$ af gericht ($Q > 0$).

\emph{2.} $k\abs{Q}/d^2$ halveren vraagt om een verdubbeld $d^2$:
$d' = 0.30\sqrt{2} \approx \qty{0.42}{m}$.
\end{solution}

\begin{solution}{exo:g11:electric-gravitational-fields:3}
$E = 600 / \num{3.0e-3} = \qty{2.0e5}{V/m}$;
$F = eE = \num{1.60e-19} \times \num{2.0e5} = \qty{3.2e-14}{N}$.
\end{solution}

\begin{solution}{exo:g11:electric-gravitational-fields:4}
$g_{\text{Mars}} = P/m = 930/250 = \qty{3.7}{N/kg}$. Astronaut:
$P = 80 \times 3.72 \approx \qty{3.0e2}{N}$ (op Aarde ongeveer
$\qty{780}{N}$).
\end{solution}

\begin{solution}{exo:g11:electric-gravitational-fields:5}
(a) Onjuist: het \omterm{def:g11:electric-gravitational-fields:field}{veld} heeft in elk punt één richting, dus snijden de
lijnen elkaar nooit. (b) Onjuist: ze vertrekken uit $+$ en komen aan bij
$-$. (c) Juist: opeenhoping codeert sterkte.
\end{solution}

\begin{solution}{exo:g11:electric-gravitational-fields:6}
\emph{1.} $d = \num{6.37e6} + \num{4.2e5} = \qty{6.79e6}{m}$, dus
$g = \num{6.67e-11} \times \num{5.97e24} / (\num{6.79e6})^2
\approx \qty{8.6}{N/kg}$ --- $88\%$ van de waarde aan het oppervlak.

\emph{2.} Ze zweven doordat station en astronaut samen vallen (vrije
val), niet doordat de zwaartekracht ontbreekt.
\end{solution}

\begin{solution}{exo:g11:electric-gravitational-fields:7}
$E = m_e g / e = \num{9.11e-31} \times 9.81 / \num{1.60e-19}
\approx \qty{5.6e-11}{V/m}$: zo'n $\num{4e14}$ keer zwakker dan het \omterm{def:g11:electric-gravitational-fields:field}{veld}
in de \omterm{def:g10:signals-and-waves:oscilloscope}{oscilloscoop}. De zwaartekracht is in de elektronenfysica volstrekt
verwaarloosbaar.
\end{solution}

\begin{solution}{exo:g11:electric-gravitational-fields:8}
\emph{1.} $E = \num{5.0e3} / \num{2.0e-2} = \qty{2.5e5}{V/m}$; in rust
geldt $qE = mg$, dus
$q = \num{2.0e-6} \times 9.81 / \num{2.5e5} \approx \qty{7.8e-11}{C}$.

\emph{2.} De \omterm{def:g10:inertia:force}{kracht} op het positieve deeltje moet omhoog wijzen, dus
wijst $\vect E$ omhoog: de \emph{onderste} plaat is de positieve.

\emph{3.} $n = \num{7.85e-11} / \num{1.60e-19} \approx \num{4.9e8}$
\omterm{def:g11:fundamental-interactions:elementary}{elementaire ladingen}.
\end{solution}

\begin{solution}{exo:g11:electric-gravitational-fields:9}
Buiten het lijnstuk wijzen beide \omterm{def:g11:electric-gravitational-fields:field}{velden} dezelfde kant op; ertussen werken
ze tegen elkaar in --- en het \omterm{def:g10:inertia:compensate}{evenwicht} vraagt om een punt dichter bij de
zwakkere bron $+Q$. Met $x$ de afstand tot $+Q$:
$kQ/x^2 = 4kQ/(0.30 - x)^2$, dus $0.30 - x = 2x$ en $x = \qty{0.10}{m}$
vanaf $+Q$.
\end{solution}

\begin{solution}{exo:g11:electric-gravitational-fields:10}
$F_e = ke^2/d^2 = \num{8.99e9} \times (\num{1.60e-19})^2 /
(\num{5.3e-11})^2 \approx \qty{8.2e-8}{N}$;
$F_g = G m_p m_e / d^2 = \num{6.67e-11} \times \num{1.67e-27} \times
\num{9.11e-31} / (\num{5.3e-11})^2 \approx \qty{3.6e-47}{N}$.
Verhouding $F_e / F_g \approx \num{2.3e39}$.
\end{solution}

\begin{solution}{exo:g11:electric-gravitational-fields:11}
$E = \sqrt{300^2 + 400^2} = \qty{500}{N/C}$, onder
$\tan^{-1}(400/300) \approx \ang{53}$ ten noorden van het oosten. Op
$q < 0$: $F = \num{2.0e-6} \times 500 = \qty{1.0e-3}{N}$, tegengesteld
aan het \omterm{def:g11:electric-gravitational-fields:field}{veld} --- $\ang{53}$ ten zuiden van het westen.
\end{solution}

\begin{solution}{exo:g11:electric-gravitational-fields:12}
$GM/(R+h)^2 = \tfrac12\,GM/R^2$ geeft $R + h = R\sqrt{2}$, dus
$h = (\sqrt{2} - 1)R \approx \qty{2.6e6}{m}$ (ongeveer \qty{2600}{km}).
Op $h = R$ is $d = 2R$: een kwart van de waarde aan het oppervlak.
\end{solution}

\begin{solution}{exo:g11:electric-gravitational-fields:13}
\emph{1.} $AM = \sqrt{3.0^2 + 4.0^2} = \qty{5.0}{cm}$; elke lading wekt
$E_1 = \num{8.99e9} \times \num{2.0e-9} / 0.050^2 \approx
\qty{7.2e3}{N/C}$ op.

\emph{2.} De componenten langs $AB$ heffen elkaar op door de symmetrie;
elk \omterm{def:g11:electric-gravitational-fields:field}{veld} draagt $E_1 \times 4/5$ langs de middelloodlijn bij, van het
lijnstuk af: $E = 2 \times \num{7.19e3} \times 0.80 \approx
\qty{1.2e4}{N/C}$, gericht langs de middelloodlijn, van $[AB]$ af.
\end{solution}

\begin{solution}{exo:g11:electric-gravitational-fields:14}
De vrije ladingen van de doos bewegen onder het uitwendige \omterm{def:g11:electric-gravitational-fields:field}{veld}: $+$
hoopt zich op de ene wand op, $-$ op de andere, en die verplaatste
ladingen wekken een inwendig \omterm{def:g11:electric-gravitational-fields:field}{veld} op dat tegen het uitwendige in gaat. Ze
blijven bewegen tot het totale \omterm{def:g11:electric-gravitational-fields:field}{veld} binnenin precies nul is --- dan is er
\omterm{def:g10:inertia:compensate}{evenwicht}. De zwaartekracht biedt maar één teken van bron: elke massa
\emph{voegt} een aantrekkend \omterm{def:g11:electric-gravitational-fields:field}{veld} toe, geen enkele kan ertegenin gaan.
Een zwaartekrachtschild zou negatieve massa nodig hebben, en die bestaat
niet.
\end{solution}

\begin{solution}{exo:g11:electric-gravitational-fields:15}
\emph{1.} Het \omterm{def:g10:universal-gravitation:weight}{gewicht} $m\vect g$ (omlaag), de \omterm{def:g10:inertia:inventory}{spankracht} $\vect T$ (langs
de draad) en de elektrische \omterm{def:g10:inertia:force}{kracht} $q\vect E$ (horizontaal). In rust
geldt $T\sin\theta = qE$ en $T\cos\theta = mg$; deel die op elkaar:
$\tan\theta = qE/(mg)$.

\emph{2.} $\tan\theta = \num{2.0e-7} \times \num{1.0e4} /
(\num{5.0e-4} \times 9.81) \approx 0.41$, dus
$\theta \approx \ang{22}$; $T = mg/\cos\theta \approx \qty{5.3e-3}{N}$.
\end{solution}

\begin{solution}{pb:g11:electric-gravitational-fields:1}
\textbf{1.} Homogeen, loodrecht op de platen, omlaag gericht (van de
positieve bovenplaat naar de negatieve onderplaat); homogeen betekent dat
de \omterm{def:g10:inertia:force}{kracht} op het druppeltje overal even groot is, waar het ook heen
drijft.
\textbf{2.} $E = U/d = 800 / \num{8.0e-3} = \qty{1.0e5}{V/m}$.
\textbf{3.} De lading is negatief, dus is de \omterm{def:g10:inertia:force}{kracht} \emph{tegengesteld}
aan $\vect E$: omhoog. Alleen met de $+$-plaat bovenaan wijst $\vect E$
omlaag en werkt de elektrische \omterm{def:g10:inertia:force}{kracht} het \omterm{def:g10:universal-gravitation:weight}{gewicht} tegen.
\textbf{4.} $P = mg = \num{4.90e-15} \times 9.81 \approx
\qty{4.8e-14}{N}$.
\textbf{5.} $\abs{q}E = mg$.
\textbf{6.} $\abs{q} = mg/E = mgd/U$.
\textbf{7.} $\abs{q} = \num{4.81e-14} \times \num{8.0e-3} / 800
\approx \qty{4.8e-19}{C}$.
\textbf{8.} $n = \num{4.8e-19} / \num{1.60e-19} = 3$ elektronen te veel.
\textbf{9.} $E$ groeit, dus $\abs{q}E > mg$: de resulterende \omterm{def:g10:inertia:force}{kracht} wijst
omhoog.
\textbf{10.} Uit $\abs{q}\,U/d = mg$ volgt $U = mgd/\abs{q}$ ---
omgekeerd evenredig met de lading.
\textbf{11.} $\abs{q} = mgd/U = \num{3.85e-16}/1200 \approx
\qty{3.2e-19}{C} = 2e$: gezakt van $3e$, dus het druppeltje heeft één
elektron \emph{verloren}.
\textbf{12.} $\abs{q} = \num{3.85e-16}/U$: \qty{3.2e-19}{C},
\qty{4.8e-19}{C}, \qty{6.4e-19}{C}, \qty{8.0e-19}{C},
\qty{9.6e-19}{C}.
\textbf{13.} Het zijn $2e, 3e, 4e, 5e, 6e$ met
$e = \qty{1.60e-19}{C}$.
\textbf{14.} Elk stapje verandert $\abs{q}$ met precies
\qty{1.6e-19}{C}: er landt telkens één elektron.
\textbf{15.} Daarvoor zou $\abs{q} = \num{3.85e-16}/700 \approx
\qty{5.5e-19}{C} = 3.4\,e$ nodig zijn --- geen geheel aantal elektronen,
dus kan geen enkele stralingspuls dat ooit voortbrengen.
\textbf{16.} Lading is \emph{gekwantiseerd}: ze komt in gehele veelvouden
van de \omterm{def:g11:fundamental-interactions:elementary}{elementaire lading} $e = \qty{1.60e-19}{C}$ --- de grootheid die
Millikan heeft gemeten.
\textbf{17.} $q \leftrightarrow m$, $E \leftrightarrow g$, en
$\abs{q}E \leftrightarrow mg$: het \omterm{def:g10:universal-gravitation:weight}{gewicht} zelf.
\textbf{18.} $g_{\text{bol}} = GM/d^2 = \num{6.67e-11} \times 1000 /
1.0^2 = \qty{6.7e-8}{N/kg}$; de \omterm{def:g10:inertia:force}{kracht} op het druppeltje is
$\num{4.90e-15} \times \num{6.7e-8} \approx \qty{3.3e-22}{N}$ ---
ongeveer $\num{1.5e8}$ keer kleiner dan zijn \omterm{def:g10:universal-gravitation:weight}{gewicht}. Hopeloos.
\textbf{19.} Omdat $G$ zo klein is, zou alleen een massa ter grootte van
een planeet boven het druppeltje de trek van de Aarde kunnen evenaren ---
en geen massa stoot af. Door de enorme $k$ en de twee tekens van lading
kan een plaat op een tafelblad de planeet in beide richtingen
overtroeven.
\textbf{20.} Eén elektron: $\Delta\abs{q}/\abs{q} = e/3e \approx 33\%$;
één oliemolecuul: $\Delta m/m = \num{5e-25} / \num{4.90e-15}
\approx \num{e-10}$. De ladingstrap verschuift de evenwichtsspanning met
honderden \omterm{def:g11:circuits-and-power:voltage}{volt}; de massatrap is bij geen enkele \omterm{def:g11:circuits-and-power:voltage}{spanning} zichtbaar. Het
\omterm{def:g10:inertia:compensate}{evenwicht} onderscheidt dus wel het \omterm{def:g11:fundamental-interactions:ladder}{atoom} van de elektriciteit ---
$e = \qty{1.60e-19}{C}$ --- terwijl de \omterm{def:g11:fundamental-interactions:ladder}{moleculen} van de olie tot een
continuüm vervagen.
\end{solution}
